\section{Solución numérica de las ecuaciones de la geodésica}

La solución analítica de las geodésicas de tipo tiempo proporciona una referencia cerrada, pero el desarrollo del trabajo requiere además integrar numéricamente las ecuaciones de movimiento \eqref{eq:geodesica_r_mino}, \eqref{eq:geodesica_theta_mino}, \eqref{eq:geodesica_t_mino} y \eqref{eq:geodesica_phi_mino}, tanto para tratar condiciones iniciales generales como para contrastar ambas ramas. La integración numérica introduce en cada paso un error de truncamiento que se acumula a lo largo de la órbita, de modo que la elección del integrador determina la fidelidad de la trayectoria a largo plazo. En esta sección seleccionamos el integrador que emplearemos como base, usando la conservación de la constante de Carter $Q$ de \eqref{eq:carter_constant} como criterio de calidad, por ser la cantidad de mayor relevancia para el análisis en la métrica de Kerr.

Para el estudio consideramos un agujero negro de Kerr con parámetro de rotación $a = 0.5$ y masa $M = 1$ en unidades geométricas. Fijamos una condición inicial en elementos osculantes con un semieje mayor igual a cuatro veces el radio de Schwarzschild, que produce una órbita que cae lentamente hacia el agujero negro. Como criterio de comparación medimos el error máximo acumulado de $Q$ en función del número de períodos orbitales, entendidos como los períodos de la órbita kepleriana asociada a la condición inicial, hasta alcanzar $10^4$ períodos. Empleamos el error máximo, y no un promedio, porque el error relativo de $Q$ fluctúa de forma natural a lo largo de la órbita, de manera que el valor pico acota la desviación efectiva de la cantidad conservada.

\subsubsection{Comparación de integradores}

Comparamos los seis integradores de propósito general disponibles en la rutina \texttt{solve\_ivp} de \texttt{SciPy} \cite{virtanen2020scipy}: los métodos explícitos de Runge-Kutta RK45, RK23 y DOP853, el método implícito de Runge-Kutta de tipo Radau, el método de diferenciación regresiva BDF y el método adaptativo LSODA. Una primera tanda de integraciones se extendió hasta $500$ períodos con una discretización de $10^4$ puntos.

Los métodos BDF, RK23 y RK45 no conservan la constante de Carter de forma adecuada: presentan fluctuaciones grandes del error entre el primer período y el centésimo. Más allá del salto abrupto en el error, gatillado por el paso físico de la órbita por la ergosfera y común a todos los métodos, el descarte de estos integradores se basa en la magnitud de esas fluctuaciones previas. Por ello los descartamos para la integración de geodésicas en la métrica de Kerr. En cambio, los métodos DOP853, Radau y LSODA conservan $Q$ de manera comparable durante los primeros cien períodos, con fluctuaciones del error de magnitud similar.

La distinción entre estos tres métodos aparece en integraciones largas, hasta $10^4$ períodos con $30\,000$ puntos. A partir de $10^3$ períodos el error del método DOP853 crece de forma pronunciada, lo que lo vuelve inadecuado para tiempos de integración prolongados. El método de Radau se mantiene como el de mejor comportamiento, sin crecimiento significativo del error acumulado, seguido por LSODA, que exhibe un crecimiento moderado. Concluimos que el método implícito de Radau es el integrador más apropiado para las geodésicas en la métrica de Kerr, y lo adoptamos como base de la integración numérica.

\subsubsection{Número de divisiones para el integrador de Radau}

Una vez fijado el método de Radau, resta determinar la discretización adecuada, es decir, el número de divisiones que equilibra el costo de la integración con la conservación de $Q$ según el horizonte temporal requerido. Para ello barremos el número de divisiones y el número de períodos orbitales, desde un período hasta $10^4$ períodos, registrando el error máximo acumulado de la constante de Carter en cada combinación.

El resultado se organiza como un mapa del error máximo acumulado en el plano definido por el número de divisiones y el número de períodos orbitales. Sobre este mapa trazamos contornos de tolerancia que, fijado un umbral de error admisible según el criterio del análisis, indican el número de divisiones mínimo requerido para cada horizonte de integración. De esta forma seleccionamos la discretización óptima en función de la duración de la órbita que se desee integrar.

El método de Radau constituye así la base de la integración numérica de las geodésicas en la métrica de Kerr. No obstante, como todo integrador de propósito general, no fue diseñado para conservar exactamente las constantes de movimiento del sistema. En la siguiente sección le acoplamos un método de proyección que fuerza al estado a permanecer sobre la subvariedad definida por las cuatro cantidades conservadas, corrigiendo la deriva residual que subsiste incluso con el mejor de los integradores evaluados.
